% Mathematical source: ../1-practice_activities/G11_T2_L1_C2_practice.tex
% User override: one worked, one guided, eight independent problems (1--10).
\makeatletter
\def\input@path{{../../../preamble/}}
\makeatother
\input{preamble.tex}
\usepackage{needspace}
\geometry{letterpaper,margin=0.25in}
\pagestyle{empty}
\linespread{1.04}
\setlength{\parskip}{4pt}
\raggedbottom
\newif\ifprintsolutions
\printsolutionsfalse
\newcommand{\stageheading}[1]{\par\addvspace{0.5\baselineskip}{\centering\large\bfseries #1\par}\nopagebreak[4]}
\newcommand{\problemtitle}[2]{\par\noindent\textbf{#1. #2}\par\nopagebreak[4]}
\newcommand{\workspace}[1]{\par\vspace{#1}\par}
\newcommand{\task}[2]{\par\textbf{Task #1 -- #2}\par}
\newcommand{\choices}[4]{\par\begingroup\small
\noindent A. #1\par B. #2\par C. #3\par D. #4\par\endgroup}
\newcommand{\infotable}[1]{\par\begingroup\small
\setlength{\tabcolsep}{4pt}\renewcommand{\arraystretch}{1.05}
\noindent\begin{tabularx}{\linewidth}{@{}p{.075\linewidth}p{.19\linewidth}p{.18\linewidth}X@{}}
\toprule\textbf{Symbol}&\textbf{Concept}&\textbf{Type}&\textbf{Meaning in context}\\\midrule
#1
\bottomrule\end{tabularx}\par\endgroup}
\newcommand{\indtable}{\infotable{
\rule{0pt}{18pt}&&Population parameters&\\\cline{1-2}\cline{4-4}
\rule{0pt}{18pt}&&&\\\midrule
\rule{0pt}{18pt}&&Sample statistics&\\\cline{1-2}\cline{4-4}
\rule{0pt}{18pt}&&&\\\cline{1-2}\cline{4-4}
\rule{0pt}{18pt}&&&\\\midrule
\rule{0pt}{18pt}&&Procedure Settings&\\\cline{1-2}\cline{4-4}
\rule{0pt}{18pt}&&&\\\cline{1-2}\cline{4-4}
\rule{0pt}{18pt}&&&\\\cline{1-2}\cline{4-4}
\rule{0pt}{18pt}&&&\\}}
\newcommand{\indtaskone}{\task{1}{Organize the information and calculate SE.}
Complete the four-column table using the convention on page 1. Identify the unknown target mean, distinguish it from the simulation mean, and show SE with substitution and units. Record shared values and level-specific settings.\indtable}
\newcommand{\indtasktwo}{\task{2}{Calculate and interpret intervals.}
For each level, show MoE, endpoints, width, and units. Interpret the range as plausible values for the unknown population mean under the model; distinguish the point estimate from the interval estimate.\space}
\newcommand{\indtaskthree}{\task{3}{Calculate observed coverage and interpret confidence.}
Show the covering count, denominator, and percentage for each procedure. Select one option and explain.\space}
\begin{document}
{\centering{\Large\bfseries C2 -- Point and interval estimation\par}
Global Economics \textbar{} Grade 11 \textbar{} Term 2 \textbar{} Lesson 1\par}
\textbf{C2:} Compare point and interval estimation and interpret confidence through long-run coverage, distinguishing this from the probability that a particular calculated interval contains the parameter.

\stageheading{Equations and basic procedure}
Use independent random observations, known population SD $\sigma$, and a normal sampling distribution of $\bar{x}$ (exact for a normal population). The target $\mu$ is fixed and unknown; $\bar{x}$ is its point estimate. All problems below use these assumptions; simulation populations are normal with the stated known fixed means and the same $\sigma$ and $n$ as their corresponding original sample.
\[
\mathrm{SE}=\frac{\sigma}{\sqrt n},\qquad
\mathrm{MoE}=z^*\mathrm{SE},\qquad
[L,U]=[\bar{x}-\mathrm{MoE},\bar{x}+\mathrm{MoE}],\qquad
\text{width}=U-L=2\mathrm{MoE}.
\]
\[
\alpha=1-C,\qquad \text{each tail}=\alpha/2,\qquad
\text{observed coverage}=\frac{\text{intervals containing the simulation mean}}{\text{intervals generated at that level}}\times100\%.
\]
\begin{center}\begin{tabular}{@{}lccc@{}}\toprule
Confidence&90\%&95\%&99\%\\\midrule
$z^*$&1.645&1.960&2.576\\\bottomrule\end{tabular}\end{center}
\textbf{Procedure.} (1) Organize information and calculate SE. (2) Calculate MoE, interval, and width; interpret and compare. (3) Use simulation evidence for observed coverage, then judge the confidence statement. A simulation interval covers when $L\leq\mu_{\rm sim}\leq U$.

\textbf{Four-column convention (problems 2--10).} Use \textbf{Symbol, Concept, Type, Meaning in context}. Give each quantity its own row in this order: $\mu,\sigma$ (Population parameters); $\bar{x},n,\mathrm{SE}$ (Sample statistics); $C,\alpha,\alpha/2,z^*$ (Procedure Settings). Write the Type only in the first row of each group and separate groups with a rule. In Meaning, explain the target population and units, record values and calculations, and distinguish shared from level-specific settings. No sample SD $s$ is provided or needed.

\textbf{Precision.} Treat the supplied $z^*$ values as the calculation values. Keep all digits unless reporting-rounding is requested. A reported width is reported upper endpoint minus reported lower endpoint. Nominal confidence and simulation coverage refer to the unrounded procedures. Use a separate denominator at each confidence level; a finite run need not equal nominal confidence.

\vspace{-3mm}
\stageheading{Worked example}
\vspace{-3mm}
\problemtitle{1}{Investigate the cycling interval}
A random sample of $n=49$ cycling trips has $\bar{x}=40$ minutes. The population SD is known: $\sigma=14$ minutes. Assume independent trips and a normal sampling model. Estimate the unknown mean cycling time $\mu$ at 95\% confidence, using $z^*=1.960$.
\[
\mathrm{SE}=\frac{14}{\sqrt{49}}=2\text{ minutes}.
\]
\[
\mathrm{MoE}=1.960(2)=3.92\text{ minutes},\qquad [40-3.92,40+3.92]=[36.08,43.92]\text{ minutes}.
\]
\textbf{Key insights.} The point estimate $\bar{x}=40$ minutes estimates the fixed unknown $\mu$; the interval gives a range of plausible values under the model. Its endpoints are equally far from 40, as required. The SD $\sigma$ describes individual trips; SE describes sample means. In the long run, about 95\% of intervals made by this procedure cover the fixed mean; this calculated interval either covers or misses it. Higher confidence widens the interval with the same centre; larger $n$ narrows it when other settings stay fixed.

\vspace{-3mm}
\stageheading{Guided Example}
\vspace{-3mm}
\problemtitle{2}{Two reports on weekly printing costs}
A random sample of $n=100$ school offices gives a mean weekly printing cost $\bar{x}=86$ thousand COP, with known population SD $\sigma=32$ thousand COP. Estimate the unknown mean for all offices in the district. Use the same sample for 90\% ($z^*=1.645$) and 99\% ($z^*=2.576$) reports. Apply the assumptions on page 1 and report endpoints to the nearest thousand COP.

\textbf{Separate simulation.} A normal population has fixed known mean $\mu_{\rm sim}=88$ thousand COP and $\sigma=32$ thousand COP. Generate 200 fresh independent samples of 100 offices. Each sample produces both unrounded intervals. The 90\% procedure misses the fixed mean 23 times; the 99\% procedure misses it 5 times.
\task{1}{Identify and organize information; calculate SE.}
Complete Meaning in context. Identify which mean is unknown; put level-specific settings in the same row.
\infotable{
$\mu$&Population mean&Population parameters&Unknown mean of \rule{1.5in}{.3pt} (units: \rule{.65in}{.3pt})\\[4pt]
$\sigma$&Population SD&&Known value: \rule{1.6in}{.3pt}\\[4pt]\midrule
$\bar{x}$&Sample mean&Sample statistics&Point estimate: \rule{1.7in}{.3pt}\\[4pt]
$n$&Sample size&&\rule{2in}{.3pt}\\[4pt]
$\mathrm{SE}$&Standard error&&$32/\sqrt{100}=$ \rule{1.1in}{.3pt} (include units)\\[4pt]\midrule
$C$&Confidence level&Procedure Settings&90\%: $0.90$; 99\%: \rule{1in}{.3pt}\\[4pt]
$\alpha$&Significance level&&90\%: $1-0.90=$ \rule{.6in}{.3pt}; 99\%: \rule{.7in}{.3pt}\\[4pt]
$\alpha/2$&Tail area&&90\%: \rule{1in}{.3pt}; 99\%: \rule{1in}{.3pt}\\[4pt]
$z^*$&Critical value&&90\%: \rule{1in}{.3pt}; 99\%: \rule{1in}{.3pt}\\[4pt]}
\task{2}{Calculate and interpret intervals; compare reporting.}
Start with $\mathrm{MoE}_{90}=1.645(\mathrm{SE})$ and $\mathrm{MoE}_{99}=2.576(\mathrm{SE})$. Show both substitutions, then complete the table (all monetary amounts in thousand COP).
\workspace{0.25in}
\begin{tabularx}{\linewidth}{@{}lXXXX@{}}\toprule
Level&MoE&Unrounded $[L,U]$&Unrounded width&Reported $[L,U]$; width\\\midrule
90\%&&&&\\[24pt]
99\%&&&&\\[24pt]\bottomrule\end{tabularx}

Explain the common centre, which interval is wider and why, how rounding affects the reported widths, and what the intervals estimate.

\task{3}{Calculate observed coverage and interpret confidence.}
90\% covering count $=200-23=$ \rule{.6in}{.3pt}; rate $=(\rule{.5in}{.3pt}/200)100\%=$ \rule{.6in}{.3pt}.
\par 99\% covering count $=200-5=$ \rule{.6in}{.3pt}; rate $=$ \rule{.8in}{.3pt}. Select and explain one option.
\choices{The 99\% interval contains 99\% of individual offices' printing costs.}
{Each finite simulation must cover at exactly its nominal confidence level.}
{The wider 99\% procedure covers the fixed mean about 99\% of the time in the long run; each calculated interval either covers or misses.}
{The reported 99\% interval has a 99\% probability of containing its fixed unknown mean.}
\workspace{0.3in}

\vspace{-3mm}
\stageheading{Independent Practice}
\problemtitle{3}{Repeated estimates of locker rental charges}
A random sample of $n=64$ customers gives $\bar{x}=72$ thousand COP for the weekly locker rental charge for all customers at a sports centre. The known population SD is $\sigma=20$ thousand COP. Estimate the unknown population mean at 95\% confidence ($z^*=1.960$), using the assumptions on page 1.
\textbf{Separate simulation.} Five fresh samples, each of size 64, come from a normal population with fixed known mean $\mu_{\rm sim}=70$ thousand COP and the same SD. Their 95\% interval endpoints are tabulated below (thousand COP). These are separate from the original sample.
\begin{center}\begin{tabular}{@{}lll@{}}\toprule
Sample & Lower endpoint & Upper endpoint\\\midrule
1 & 64.1 & 73.9\\
2 & 67.1 & 76.9\\
3 & 70.1 & 79.9\\
4 & 63.1 & 72.9\\
5 & 65.1 & 74.9\\
\bottomrule\end{tabular}\end{center}
\indtaskone
\indtasktwo

\indtaskthree
Count the covering intervals in the table. Which statement is correct?
\choices{Four of five cover; 80\% is therefore the long-run coverage.}{Four of five cover; the procedure covers about 95\% in the long run, and this particular interval either covers or misses.}{Four of five cover; the original interval has a 95\% probability of containing its fixed mean.}{All five cover; the procedure covers about 95\% in the long run.}
\workspace{0.35in}

\clearpage
\problemtitle{4}{Repeated estimates of nursery delivery charges}
A random sample of $n=81$ orders gives $\bar{x}=54$ thousand COP for the delivery charge for all nursery orders in a specified month. The known population SD is $\sigma=18$ thousand COP. Estimate the unknown population mean at 95\% confidence ($z^*=1.960$), using the assumptions on page 1.
\textbf{Separate simulation.} Five fresh samples, each of size 81, come from a normal population with fixed known mean $\mu_{\rm sim}=55$ thousand COP and the same SD. Their 95\% interval endpoints are tabulated below (thousand COP). These are separate from the original sample.
\begin{center}\begin{tabular}{@{}lll@{}}\toprule
Sample & Lower endpoint & Upper endpoint\\\midrule
1 & 50.08 & 57.92\\
2 & 51.08 & 58.92\\
3 & 46.08 & 53.92\\
4 & 52.08 & 59.92\\
5 & 55.08 & 62.92\\
\bottomrule\end{tabular}\end{center}
\indtaskone
\indtasktwo
\workspace{1.05in}
\indtaskthree
Count the covering intervals in the table. Which statement is correct?
\choices{Three of five cover; the calculated interval has a 95\% probability of containing the fixed mean.}{Five of five cover; confidence is about repeated intervals.}{Three of five cover; the procedure covers about 95\% in the long run, while a fixed calculated interval either covers or misses.}{Three of five cover; 60\% replaces 95\% as the long-run confidence.}
\workspace{0.35in}

\clearpage
\problemtitle{5}{Two reports on monthly pet-care spending}
A random sample of $n=144$ households has $\bar{x}=145$ thousand COP and known population SD $\sigma=42$ thousand COP. Estimate the unknown mean monthly pet-care spending among all registered households in a district. Both reports use this same sample: 90\% ($z^*=1.645$) and 99\% ($z^*=2.576$). Use the assumptions on page 1. Round reported endpoints to the nearest thousand COP.
\textbf{Separate simulation.} A normal population has fixed known mean $\mu_{\rm sim}=148$ thousand COP and the same SD. From 150 fresh samples of size 144, construct both unrounded intervals for each sample. The 90\% intervals miss the fixed mean 17 times; the 99\% intervals miss it 3 times. Check coverage before rounding and separately at each level.
\indtaskone
\indtasktwo
Also state the unrounded endpoints and widths, reported endpoints and reported widths. Compare both widths before and after rounding and explain the common centre.
\workspace{1.55in}
\indtaskthree
Which interpretation is correct?
\choices{The wider 99\% interval contains 99\% of individual household spending amounts.}{Both finite coverage rates must equal their nominal confidence levels.}{The wider reported interval has a 99\% probability of containing the fixed mean.}{The 99\% unrounded procedure covers the fixed mean about 99\% of the time in the long run; rounding is a separate reporting step.}
\workspace{0.35in}

\clearpage
\problemtitle{6}{Two reports on weekly bakery supply costs}
A random sample of $n=64$ bakeries has $\bar{x}=230$ thousand COP and known population SD $\sigma=36$ thousand COP. Estimate the unknown mean weekly supply cost for all neighbourhood bakeries in a city. Both reports use this same sample: 90\% ($z^*=1.645$) and 99\% ($z^*=2.576$). Use the assumptions on page 1. Round reported endpoints to the nearest thousand COP.
\textbf{Separate simulation.} A normal population has fixed known mean $\mu_{\rm sim}=228$ thousand COP and the same SD. From 250 fresh samples of size 64, construct both unrounded intervals for each sample. The 90\% intervals miss the fixed mean 28 times; the 99\% intervals miss it 4 times. Check coverage before rounding and separately at each level.
\indtaskone
\indtasktwo
Also state the unrounded endpoints and widths, reported endpoints and reported widths. Compare both widths before and after rounding and explain the common centre.
\workspace{1.55in}
\indtaskthree
Which interpretation is correct?
\choices{The 99\% unrounded procedure has greater long-run coverage and gives a wider interval from this sample; each calculated interval either covers or misses.}{The 99\% interval contains 99\% of individual bakeries' costs.}{The simulation must give exactly 90\% and 99\% coverage.}{The reported 99\% interval has a 99\% probability of containing its fixed mean.}
\workspace{0.35in}

\clearpage
\problemtitle{7}{Clinic waiting times in a simulation}
A random sample of $n=49$ appointments gives $\bar{x}=26$ minutes; known population SD is $\sigma=14$ minutes. Estimate the unknown mean waiting time for all appointments at a clinic in a specified month. Use 90\% confidence ($z^*=1.645$) and the assumptions on page 1.
\textbf{Separate simulation.} Twenty fresh samples of size 49 come from a normal population with fixed known mean $\mu_{\rm sim}=27$ minutes and the same SD. The table gives counts of unrounded 90\% intervals containing that mean, grouped by generation order.
\begin{center}\begin{tabular}{@{}lllll@{}}\toprule
Samples in order & 1--5 & 6--10 & 11--15 & 16--20\\\midrule
Covering intervals & 4 & 5 & 4 & 3\\
\bottomrule\end{tabular}\end{center}
\indtaskone
\indtasktwo
\workspace{1.35in}
\indtaskthree
Total the grouped counts. Assess the claim: ``Every run of 20 must have exactly 18 covering intervals. Which option is correct?
\choices{Eighteen of 20 cover; every 90\% run must give exactly 18.}{Sixteen of 20 cover; finite runs may differ, while long-run coverage is about 90\% under the model.}{Sixteen of 20 cover; the observed 80\% is the long-run coverage.}{Sixteen of 20 cover; each fixed interval has a 90\% chance of containing the fixed mean.}
\workspace{0.35in}

\clearpage
\problemtitle{8}{Parcel sorting times in a simulation}
A random sample of $n=64$ parcels gives $\bar{x}=18$ minutes; known population SD is $\sigma=12$ minutes. Estimate the unknown mean sorting time for all parcels at a depot during a specified shift. Use 90\% confidence ($z^*=1.645$) and the assumptions on page 1.
\textbf{Separate simulation.} Twenty fresh samples of size 64 come from a normal population with fixed known mean $\mu_{\rm sim}=19$ minutes and the same SD. The table gives counts of unrounded 90\% intervals containing that mean, grouped by generation order.
\begin{center}\begin{tabular}{@{}lllll@{}}\toprule
Samples in order & 1--5 & 6--10 & 11--15 & 16--20\\\midrule
Covering intervals & 5 & 4 & 5 & 5\\
\bottomrule\end{tabular}\end{center}
\indtaskone
\indtasktwo
\workspace{1.35in}
\indtaskthree
Total the grouped counts. Assess the claim: ``Every run of 20 must have exactly 18 covering intervals. Which option is correct?
\choices{Nineteen of 20 cover; every 90\% run must give exactly 18, so the procedure is invalid.}{Eighteen of 20 cover; each run must match 90\%.}{Nineteen of 20 cover; 95\% becomes the procedure's long-run confidence.}{Nineteen of 20 cover; the count may vary across finite runs even though long-run coverage is about 90\%.}
\workspace{0.35in}

\clearpage
\problemtitle{9}{Three confidence levels for library processing times}
A random sample of $n=100$ books has $\bar{x}=24$ minutes and known population SD $\sigma=15$ minutes. Estimate the unknown mean processing time for all returned books at a library during one month. Construct 90\%, 95\%, and 99\% intervals from this same sample, using $z^*=1.645$, $1.960$, and $2.576$, respectively, and the assumptions on page 1. Keep all calculated endpoint digits.
\textbf{Separate simulation.} From a normal population with fixed known mean $\mu_{\rm sim}=25$ minutes and the same SD, generate 250 fresh samples of size 100. Construct all three unrounded intervals from each fresh sample, centred on that sample\textquotesingle s own mean. The categories below are mutually exclusive; each sample appears once.
\begin{center}\begin{tabular}{@{}ll@{}}\toprule
Intervals containing the fixed simulation mean & Samples\\\midrule
All three & 219\\
Only 95\% and 99\% & 18\\
Only 99\% & 10\\
None & 3\\
\bottomrule\end{tabular}\end{center}
\indtaskone
\indtasktwo
Compare the three widths and centres; explain their nesting.
\workspace{0.9in}
\indtaskthree
Derive all three covering counts. Which option gives the rates (90\%, 95\%, 99\%) and a correct interpretation? Round percentages to two decimals.
\choices{90\%, 95\%, 99\%; every finite run must match its nominal levels.}{87.6\%, 94.8\%, 98.8\%; higher confidence moves the centre closer to the unknown mean.}{87.6\%, 94.8\%, 98.8\%; higher confidence widens intervals about the same sample mean, and describes long-run coverage of the fixed mean.}{87.6\%, 94.8\%, 98.8\%; the calculated 99\% interval has a 99\% probability of containing its fixed mean.}
\workspace{0.35in}

\clearpage
\problemtitle{10}{Three confidence levels for equipment setup times}
A random sample of $n=81$ production runs has $\bar{x}=62$ minutes and known population SD $\sigma=27$ minutes. Estimate the unknown mean setup time for all production runs at a workshop during one month. Construct 90\%, 95\%, and 99\% intervals from this same sample, using $z^*=1.645$, $1.960$, and $2.576$, respectively, and the assumptions on page 1. Keep all calculated endpoint digits.
\textbf{Separate simulation.} From a normal population with fixed known mean $\mu_{\rm sim}=60$ minutes and the same SD, generate 300 fresh samples of size 81. Construct all three unrounded intervals from each fresh sample, centred on that sample\textquotesingle s own mean. The categories below are mutually exclusive; each sample appears once.
\begin{center}\begin{tabular}{@{}ll@{}}\toprule
Intervals containing the fixed simulation mean & Samples\\\midrule
All three & 268\\
Only 95\% and 99\% & 16\\
Only 99\% & 12\\
None & 4\\
\bottomrule\end{tabular}\end{center}
\indtaskone
\indtasktwo
Compare the three widths and centres; explain their nesting.
\workspace{0.9in}
\indtaskthree
Derive all three covering counts. Which option gives the rates (90\%, 95\%, 99\%) and a correct interpretation? Round percentages to two decimals.
\choices{About 89.33\%, 94.67\%, 98.67\%; confidence describes long-run coverage, and greater confidence widens intervals with the same centre.}{90\%, 95\%, 99\%; finite runs must match all three levels.}{About 89.33\%, 94.67\%, 98.67\%; higher confidence changes the sample mean to approach the fixed mean.}{About 89.33\%, 94.67\%, 98.67\%; 99\% of individual setup times lie in the calculated 99\% interval.}
\workspace{0.35in}
\end{document}
